% === Dokumentklasse ==========================================================
\documentclass[%
    a4paper,							% A4
    12pt,								% Schriftgröße 12
    twoside,							% zweiseitig
    fleqn								% Formeln linksbündig ausrichten
]{book}


% === Grundlegende Pakete =====================================================
\usepackage{etoolbox}					% portiert viele nützliche Sachen aus e-TeX (z.B. booleans, ifs)
\usepackage[							% erweiterte Angabe von Farben
    pdftex,								% Farbtreiber auswählen
    dvipsnames							% vordefinierte Farben laden
]{xcolor}								% erweiterte Angabe von Farben
\usepackage{xparse}						% high-level Interface für Dokumentbefehle wie \NewDocument(Command,Environment)

\usepackage[utf8]{inputenc}				% Kodierung der *.tex Dateien ist UTF-8 nicht ISO-8859-1
\usepackage{cmap}						% character map tables -> pdf Inhalt wird besser durchsuch- und kopierbar
\usepackage[T1]{fontenc}				% OT1 encoding für deutsche Sonderzeichen -> pdf Inhalt kann Sonderzeichen durchsucht werden
\usepackage{lmodern}					% Latin-Modern-Schriftart (Computer Modern in Verbindung mit OT1 führt zu Bitmap-Fonts auf Windows)
\usepackage{microtype}					% mikrotypografische Einstellungen, die den gesetzten Text nochmals wesentlich verbessern (weniger badboxes)


% === Benutzerdefinierte Einstellungen ========================================
% === Art der Arbeit ==========================================================
% von den nachfolgenden Blöcken bitte den richtigen auswählen und die anderen auskommentieren/löschen

% --- Diplom ----------
%\newcommand{\settingsDegree}{{Diplom}}
%\newcommand{\settingsDegreeName}{{Diplominformatiker}}
%\newcommand{\settingsDegreeName}{{Diplomingenieur}}

% --- Bachelor ----------
% \newcommand{\settingsDegree}{{Bachelor}}
% \newcommand{\settingsDegreeName}{{Bachelor of Science}}

% --- Master ----------
\newcommand{\settingsDegree}{{Master}}
\newcommand{\settingsDegreeName}{{Master of Science}}

% === Name, Abgabedatum und Sprache der Arbeit ================================
\newcommand{\settingsName}{{Pierre Helbing}}
\newcommand{\settingsFinishDate}{{DD.MM.YYYY}}
\newcommand{\settingsLanguage}{german}   		% german / american

% === Weitere Einstellungen ===================================================
% --- Suchpfad (Unterverzeichnis) für eingebundene Grafiken ----------
\newcommand{\settingsGraphicsPath}{image/}

% --- Hinweiskapitel ----------
\newbool{settingsWithHints}
\setbool{settingsWithHints}{false}				% true / false

% --- Zeilennummern ----------
\newbool{settingsWithLineNumbers}
\setbool{settingsWithLineNumbers}{false}			% true / false

% --- Todos ----------
\newbool{settingsWithTodos}
\setbool{settingsWithTodos}{true}				% true / false

% --- Anzahl an Nummerierungsebenen im Text und Inhaltsverzeichnis ----------
% 1: \section
% 2: \section + \subsection
% Achtung: 3 oder 4 nur nach Absprache mit Betreuer !
% 3: \section + \subsection + \subsubsection
% 4: \section + \subsection + \subsubsection + \paragraph
\setcounter{secnumdepth}{2}
\setcounter{tocdepth}{2}

% --- Anforderungen ----------

\usepackage{tabularx}
\usepackage{multicol}
\newenvironment{myreq}[1]{%
    \setlist[description]{font=\normalfont\color{darkgray}}%
    \begin{tcolorbox}[colframe=black,colback=white, sharp corners, boxrule=1pt]%
        \bfseries\color{blue}%
        \begin{description}#1}%
            {\end{description}\end{tcolorbox}}

\newcommand{\threeinline}[3]{\begin{multicols}{3}#1 #2 #3\end{multicols}}
\newcommand{\twoinline}[2]{\begin{multicols}{2}#1 #2\end{multicols}}

\newcommand{\reqno}{\item[Requirement \#:]}
\newcommand{\reqtype}{\item[Requirement Type:]}
\newcommand{\reqevent}{\item[Event/BUC/PUC \#:]}
\newcommand{\reqdesc}{\item[Description:]}
\newcommand{\reqrat}{\item[Rationale:]}
\newcommand{\reqorig}{\item[Originator:]}
\newcommand{\reqfit}{\item[Fit Criterion:]}
\newcommand{\reqsatis}{\item[Customer Satisfaction:]}
\newcommand{\reqdissat}{\item[Customer Dissatisfaction:]}
\newcommand{\reqdep}{\item[Dependencies:]}
\newcommand{\reqconf}{\item[Conflicts:]}
\newcommand{\reqmater}{\item[Materials:]}
\newcommand{\reqhist}{\item[History:]}
\usepackage{multicol}
\usepackage{multirow}
\usepackage[dvipsnames]{xcolor}
\usepackage{enumitem}
\usepackage{tcolorbox}

% === Übersetzungen ===========================================================
% Definitionen je nach \mylanguage (siehe NIKR_settings.tex)
\ifdefstring{\settingsLanguage}{german}{%
    \newcommand{\acroname}{Abkürzungsverzeichnis}	% Name für Abkürzungsverzeichnis
    \newcommand{\todoname}{Todo Liste}	% Name für Todo-Liste
    \newcommand{\pagename}{Seite}		% Name für Seite
}%
{%
    \newcommand{\acroname}{Acronyms}	% Name für Abkürzungsverzeichnis
    \newcommand{\todoname}{Todo list}	% Name für Todo-Liste
    \newcommand{\pagename}{page}		% Name für Seite
}%


% === Wichtige Pakete und Einstellungen =======================================
\usepackage{calc}						% ermöglicht Arithmetik in den Argumenten von Befehlen

% Einstellungen je nach \mylanguage (siehe NIKR_settings.tex)
\ifdefstring{\settingsLanguage}{german}{%
    \usepackage[ngerman]{babel}			% Spracheinstellungen (für deutsch z.B. Contents -> Inhaltsverzeichnis, etc.)
    \usepackage{bibgerm}            	% Stylefile für deutsche Literaturstellenangabe
    \bibliographystyle{gerapali}		% Stil für Literaturangaben festlegen
}%
{%
    \usepackage[american]{babel}		% Spracheinstellungen (für deutsch z.B. Contents -> Inhaltsverzeichnis, etc.)
    \bibliographystyle{apalike}			% Stil für Literaturangaben festlegen
    \frenchspacing						% einfaches Leerzeiches nach Satzende (für deutsch bereits Standard)
}%
\usepackage[%
    noadjust							% noadjust verhindert automatische Leerzeichen um die Referenz, was am Zeilenanfang zu Problemen führt
]{cite}     							% erlaubt Zeilenumbruch innerhalb von Zitierungen

\usepackage{graphicx}					% erweiterte Argumente in \in­clude­graph­ics
\graphicspath{{\settingsGraphicsPath}}	% Standard-Pfad für Bilder siehe NIKR_settings.tex

\usepackage[bf]{caption}				% erlaubt erweiterte Formatierungen in \caption (siehe unten)
\usepackage{subcaption}					% mehrere Abbildungen nebeneinander

\usepackage{amsmath}					% ermöglicht \DeclareMathOperator
\usepackage{amssymb}					% mathmatische Symbole und Sonderzeichen
\usepackage{nicefrac}					% für \nicefrac
\usepackage{nccmath}					% für \mfrac

\usepackage{icomma}						% intelligentes Komma (macht Verwendung von {,} überflüssig)
\usepackage{siunitx}					% für einheitliche Angabe von Einheiten

\usepackage{fancyhdr}					% Kopf- und Fußzeilen (siehe unten)

\usepackage{hhline}						% erweitere Rahmengestaltung in Tabellen

\usepackage[%							% Einbettung von Links im Dokument und erlaubt die Nutzugn von \url
    hyperfootnotes=false,				% keine Fußnoten als Link im Dokument (geht nicht mit footmisc)
    pagebackref=true					% Backrefs im Literaturverzeichnis
]{hyperref}           					% Einbettung von Links im Dokument und erlaubt die Nutzugn von \url
\renewcommand*{\backref}[1]{\textit{(\pagename:~#1)}}   % Format für backrefs

\usepackage[%							% erlaubt erweiterte Formatierungen von Fußnoten (siehe unten)
    multiple,							% mehrere mit Komma abtrennen
    hang								% linksbündig, \footnotemargin entscheidet über Einrückung
]{footmisc}								% erlaubt erweiterte Formatierungen von Fußnoten (siehe unten)
\patchcmd{\footref}{\ref}{\ref*}{}{}	% Hyperlink in \footref entfernen

\usepackage[nohyperlinks]{acronym}		% Abkürzungsverzeichnis

\usepackage{setspace}					% für \setstretch (ändern Zeilenabstand, aber nicht floating Umgebungen)

\usepackage[%							% Todos
    colorinlistoftodos,					% farbige Markierungen in Todo-Liste
    prependcaption,						% caption=val
    textsize=tiny,						% Schriftgröße
    linecolor=red,						% Standard-Linienfarbe für \todo
    backgroundcolor=red!25,				% Standard-Hintergrundfarbe für \todo
    bordercolor=red,					% Standard-Rahmenfarbe für \todo
    textwidth=2cm,						% Standard-Breite für \todo
]{todonotes}							% Todos
\ifbool{settingsWithTodos}{%
    \setlength{\marginparwidth}{2cm}	% sonst werden Notes am Rand nicht richtig angezeigt
    \NewDocumentCommand{\todoaddref}{O{} m}{%
        \todo[linecolor=blue,backgroundcolor=blue!25,bordercolor=blue,#1]{#2}%
    }
    \NewDocumentCommand{\todouncertain}{O{} m}{%
        \todo[linecolor=green,backgroundcolor=green!25,bordercolor=green,#1]{#2}%
    }
    \NewDocumentCommand{\todooptional}{O{} m}{%
        \todo[linecolor=cyan,backgroundcolor=cyan!25,bordercolor=cyan,#1]{#2}%
    }
    \pretocmd{\mainmatter}{\listoftodos[\todoname]{\markboth{\MakeUppercase{\todoname}}{\MakeUppercase{\todoname}}}}{}{}
}{
    \presetkeys{todonotes}{disable}{}	% disable \todo
    \NewDocumentCommand{\todoaddref}{O{} m}{}
    \NewDocumentCommand{\todouncertain}{O{} m}{}
    \NewDocumentCommand{\todooptional}{O{} m}{}
}

\usepackage[switch*,pagewise]{lineno}	% Zeilennummern
\ifbool{settingsWithLineNumbers}{%
    \renewcommand\linenumberfont{\textbf\sffamily\color{black!50}\footnotesize}
    \apptocmd{\mainmatter}{\linenumbers}{}{}
    \pretocmd{\backmatter}{\nolinenumbers}{}{}
}{}

\usepackage{placeins}					% FloatBarriers

\usepackage{enumitem}					% erweiterte Formatierung von \enumerate, \itemize und \description

\usepackage[linewidth=0.5pt]{mdframed}	% für Boxen in Hinweisen

\usepackage[ddmmyyyy]{datetime}			% Datumsangabe
\renewcommand{\dateseparator}{.}		% Punkt als Trennzeichen in Datumsangabe


% === Längen und Abstände =====================================================
% horizontales Layout
\setlength{\oddsidemargin}{0.2in}
\setlength{\evensidemargin}{0.0in}
\setlength{\textwidth}{\paperwidth - 2.2in}

% vertikales Layout
%\setlength{\topskip}{0.0cm}
\setlength{\headheight}{15.1pt}
%\setlength{\headsep}{0.0cm}
\setlength{\topmargin}{0.0cm}
\setlength{\footskip}{0.6in}
\setlength{\textheight}{\paperheight - 2.0in}
\addtolength{\textheight}{-1.0\headheight}
\addtolength{\textheight}{-1.0\headsep}
\addtolength{\textheight}{-1.0\footskip}

% Zeilenabstand
\setstretch{1.3}
\AtBeginEnvironment{tabular}{\setstretch{1.3}}

% Einrückung von Formeln
\setlength{\mathindent}{1.0cm}

% Absätze
\setlength{\parindent}{0.0cm}

% Fußnoten
\renewcommand{\footnotelayout}{\setstretch{1.2}}
\setlength\footnotemargin{10pt}

% Listen (noitemsep, nosep, ...)
\setlist{noitemsep}

\setlength{\topsep}{0.3cm}


% === Bild- und Tabellenunterschrift ==========================================
\renewcommand{\captionfont}{\small \setstretch{1.3}}
\newcommand{\NIcaption}[2]{\caption[#1]{#1\protect\\ \emph{#2}}}
\setcaptionmargin{0.75cm}


% === Abkürzungsverzeichnis ===================================================
% Verwendung vor jedem Kapitel zurücksetzen
\pretocmd{\chapter}{\acresetall}{}{}


% === Seitenstil ==============================================================
% Pagestyle plain überschreiben
\pagestyle{fancy}
% Kapitel- und Abschnittangaben ohne Punkt
\renewcommand{\sectionmark}[1]{\markright{\uppercase{\thesection~~#1}}}
\renewcommand{\chaptermark}[1]{\markboth{\uppercase{\chaptername\ \thechapter~~#1}}{}}
\fancypagestyle{plain}{%
    \fancyhead[ER]{\itshape\leftmark}%
    \fancyhead[OL]{\itshape\rightmark}%
    \fancyhead[EL,OR]{\thepage}%
    \fancyfoot[EL,OL]{}%
    \fancyfoot[EC,OC]{}%
    \renewcommand{\headrulewidth}{0.4pt}%
    \renewcommand{\footrulewidth}{0.4pt}%
}


% === Verweise ================================================================
% Klammern in Formel-Referenzen entfernen
\makeatletter
\renewcommand\tagform@[1]{\maketag@@@{\ignorespaces#1\unskip\@@italiccorr}}
\makeatother


% === Angabe von Einheiten ====================================================
\ifdefstring{\settingsLanguage}{german}{%
    \sisetup{locale=DE}		% deutsche Lokalisierung (konvertiert 1.00 automatisch zu 1,00)
}%
{%
    \sisetup{locale=US}		% englische Lokalisierung (konvertiert 1,00 automatisch zu 1.00)
}%


% === Mathematische Definitionen ==============================================
% Darstellung von Vektoren und Matrizen
\renewcommand{\vec}[1]{\underline{\mathbf{\MakeLowercase{#1}}}}		% Vektoren
\newcommand{\veci}[1]{\underline{\MakeLowercase{#1}}}				% Vektoren als Indizes
\newcommand{\mat}[1]{\underline{\mathbf{\MakeUppercase{#1}}}}		% Matrizen
\newcommand{\mati}[1]{\underline{\MakeUppercase{#1}}}				% Matrizen als Indizes
% zusätzliche mathematische Operatoren (damit sie nicht als Formelzeichen geschrieben werden)
\DeclareMathOperator{\step}{step}									% Stufenfunktion
\DeclareMathOperator{\sign}{sign}									% Vorzeichen


% === Manuelle Definition von Worttrennungen ==================================
\hyphenation{
    Convolutional
    Neural
    Net-work
    Net-works
    Ko-ef-fi-zi-ent
    Ko-ef-fi-zi-ent-en
    Drop-out
    pixel-genaue
    Patch
    Patch-größen
}


% === Pseudocode-Darstellung ==================================================
% Import nicht oben, weil \parindent zuvor gesetzt werden muss!
% siehe: https://ctan.org/pkg/algorithm2e?lang=de
\ifdefstring{\settingsLanguage}{german}{%
    \usepackage[%						% Pseudocode
        linesnumbered,					% mit Zeilennummern
        noend,							% Ende von Befehlen, wie etwa While, unterdrücken
        ruled,							% Layout
        german,							% deutsche Bezeichnung und deutsches Verzeichnis
        %onelanguage,					% Keywords übersetzen
        algochapter						% Nummerierung analog zu Abbildungen
    ]{algorithm2e}						% Pseudocode
}{%
    \usepackage[%						% Pseudocode
        linesnumbered,					% mit Zeilennummern
        noend,							% Ende von Befehlen, wie etwa While, unterdrücken
        ruled,							% Layout
        algochapter						% Nummerierung analog zu Abbildungen
    ]{algorithm2e}						% Pseudocode
}%

\newenvironment{NIalgorithm}{%
    % Algorithmus um 1.5em einrücken, damit Zeilennummern nicht im Rand sind
    \setlength{\algomargin}{1.5em}%
    % Padding oben und unten für Caption
    \setlength{\interspacetitleruled}{\smallskipamount}%
    % Padding oben und unten für Algorithmus
    \SetAlgoInsideSkip{smallskip}%
    % Kommentarstyle ändern
    \newcommand\NIcommentstyle[1]{\ttfamily\textcolor{black!60}{##1}}
    \SetCommentSty{NIcommentstyle}%
    % zweiten Teil in NIcaption unterdrücken (falls NIcaption genutzt wird)
    \renewcommand{\NIcaption}[2]{\caption[##1]{##1}}%
    \begin{algorithm}%
        % Zeilenabstand minimal vergrößern
        \setstretch{1.1}%
        % kleine Schrift
        \small%
        % Semikolon unterdrücken
        \DontPrintSemicolon%
        % korrekt ausgerichtete mehrzeilige Input- bzw. Outputdefinitionen mittels \Input und \Output
        \SetKwInOut{Input}{Input}%
        \SetKwInOut{Output}{Output}%
        % Abschnitt für In- und Outputs zurückrücken
        \pretocmd{\Input}{\Indentp{-1.5em}}{}{}%
        \apptocmd{\Output}{\Indentp{1.5em}}{}{}%
        }{%
    \end{algorithm}%
}

\begin{document}
% --- Titelseite, Danksagung, Einverständniserklärung --------------------------
\pagestyle{empty}
% Buchstaben für Seitennummerierung verwenden 
% (verhindert "destination with the same identifier (name{page.X})" Warnung)
\pagenumbering{alph}

% === Titelblatt ==============================================================
\begin{titlepage}
	\hspace{0.2cm}
	\begin{minipage}{3.5cm}
		\includegraphics[width=0.8\textwidth]{images/logo}
	\end{minipage}
	\hspace{0.2cm}
	\begin{minipage}{11cm}
		\vspace{0.7cm}
		\large
		{\bf Technische Universität Ilmenau}\newline
		Fakultät für Informatik und Automatisierung\newline
		Fachgebiet Neuroinformatik und Kognitive Robotik
	\end{minipage}
	\begin{center}
		\vspace{0.8cm}
		{\Large\bfseries Entwicklung einer CrossLab-kompatiblen integrierten Entwicklungsumgebung für das GOLDi-Remotelab\\}
		\vspace{0.8cm}
		\settingsDegree arbeit zur Erlangung des akademischen Grades \settingsDegreeName\\[0.5cm]
		{\Large \bfseries \settingsName\\[1.0cm]}
		\begin{table}[ht]
			\centering
			\begin{tabular}{ll}
				Betreuer: & Dr. Detlef Streitferdt                    \\[2mm]
				\multicolumn{2}{l}{Verantwortlicher Hochschullehrer:} \\
				          & Prof. Dr.-Ing. habil. Daniel Ziener       \\[2cm]
				\multicolumn{2}{p{13cm}}{Die \settingsDegree arbeit wurde am \settingsFinishDate \ bei der Fakultät für Informatik und Automatisierung der Technischen Universität Ilmenau eingereicht.}
			\end{tabular}
		\end{table}
		% Hinweis für Entwurfsversion ausgeben (Variablen aus NIKR_settings.tex prüfen)
		\ifboolexpr{bool{settingsWithTodos} or bool{settingsWithLineNumbers} or bool{settingsWithHints}}{%
		{\color{red} Entwurf: \today \\[2mm]
		\textbf{Dies ist nicht die finale Druckvorlage.}\\[2mm]
		{\small%
		\ifbool{settingsWithHints}{%
			Hinweiskapitel aktiviert (siehe \texttt{settingsWithHints} in \texttt{NIKR\_settings.tex})\\%
		}{}%					
		\ifbool{settingsWithLineNumbers}{%
			Zeilennummern aktiviert (siehe \texttt{settingsWithLineNumbers} in \texttt{NIKR\_settings.tex})\\%
		}{}%
		\ifbool{settingsWithTodos}{%
			Todo-Markierungen aktiviert (siehe \texttt{settingsWithTodos} in \texttt{NIKR\_settings.tex})\\%
		}{}%
		}%
		}%
		}{}
	\end{center}
\end{titlepage}

\cleardoublepage

% === Danksagung ==============================================================
% Sollten Sie diesen Abschnitt nicht nutzen wollen, kommentieren Sie alle Zeilen
% bis zur Einverstaendniserklaerung (inkl. \cleardoublepage) aus
\vspace*{5cm}

Danksagung

Dieser Abschnitt {\bf kann} genutzt werden, um denjenigen Personen Dank auszusprechen, die Sie bei der Erstellung der Arbeit unterstützt haben.

\cleardoublepage

% === Einverstaendniserklaerung ================================================
\vspace*{16cm}

\begin{tabular}{lp{12.5cm}}
	{Erklärung:} & {"`Hiermit versichere ich, dass ich diese wissenschaftliche Arbeit selbständig verfasst und nur die angegebenen

			
			
			
			
			
			
			
			
			
			
			Quellen und Hilfsmittel verwendet habe. Alle von mir aus anderen
			Veröffentlichungen übernommenen Passagen sind als solche gekennzeichnet."'}
\end{tabular}
\vspace*{1.5cm}

\begin{tabular}{l}
	Ilmenau, \settingsFinishDate \\
	\\
\end{tabular}
\hfill
\begin{tabular}{c}
	{\makebox[6.0cm]{\dotfill}} \\
	\settingsName               \\
\end{tabular}

\cleardoublepage


% --- Inhalts- und Abkürzungsverzeichnis ---------------------------------------
\frontmatter
\pagestyle{plain}
% Inhaltsverzeichnis
\tableofcontents
% Abkürzungsverzeichnis
\chapter*{\acroname}
\markboth{\MakeUppercase{\acroname}}{\MakeUppercase{\acroname}}
% 1. Definition von Abkürzungen
\begin{acronym}[XXXXXX] % durch XXXXXX kann der Einzug bestimmt wird
	\setlength{\itemsep}{-\parsep}
	\acro{IDE}{Integrated Development Environment}
	\acro{LSP}{Language Server Protocol}
	\acro{DAP}{Debug Adapter Protocol}
	\acro{VSCode}{Visual Studio Code}
	\acro{OT}{Operational Transformation}
	\acro{CRDTs}{Conflict-free Replicated Data Types}
	\acro{LMS}{Lernmanagementsysteme}
	\acro{LTI}{Learning Tools Interoperability}
\end{acronym}

% 2. Beispiel für Verwendung im Fließtext, z.B.: \ac{CNN}



% --- Inhalt -------------------------------------------------------------------
\mainmatter

% Hinweiskapitel
\ifbool{settingsWithHints}{\input{content/hinweise}}{}

% Kapitel
\chapter{Einleitung}\label{section:einleitung}

\begin{note}
    \textbf{Notizen:}
    \begin{itemize}
        \item Remote Labore (z.B. GOLDi)
        \item IDEs in Remote Laboren (z.B. GOLDi und WIDE)
        \item Änderungen und Möglichkeiten durch CrossLab
        \item Weiterer Verlauf der Arbeit
    \end{itemize}
\end{note}
\todo{Einleitung}

% Das Grid of Online Laboratory Devices Ilmenau (GOLDi) ist ein hybrides Online Labor. Es ermöglicht Studierenden reale sowie simulierte elektromechanische Hardwaremodelle mit verschiedenen Kontrolleinheiten zu steuern. Vor dem CrossLab Projekt waren nur Experiment mit einem Hardwaremodell und einer Kontrolleinheit möglich. Aufgrund der neuen CrossLab Architektur ist es allerdings nun möglich Experimente freier zu konfigurieren. Dabei werden sogenannte Laborgeräte definiert, die gewisse Services anbieten. Über diese Services können verschiedene Laborgeräte zu einem Experiment zusammengefügt werden. Nehmen wir zum Beispiel ein 3-Achs-Portal sowie einen Microcontroller. Beide diese Geräte können einen Service anbieten, der ihre Pins definiert. Über diese Services können dann die Pins der beiden Geräte verbunden werden.

% WIDE ist eine eigens für das GOLDi Remotelab entwickelte integrierte Entwicklungsumgebung (IDE). Sie ermöglicht es Nutzern direkt im Browser Programme für die Kontrolleinheiten zu schreiben. Weiterhin wird auch die Kompilierung des Quellcodes sowie das Hochladen des Kompilats auf die Kontrolleinheit unterstützt. Allerdings ist WIDE auf die alte Architektur des GOLDi Remotelab ausgelegt, in welcher ein Experiment nur aus jeweils einem Hardwaremodell und einer Kontrolleinheit bestand. In der neuen CrossLab Architektur kann jedoch ein Experiment mehrere Hardwaremodelle und Kontrolleinheiten enthalten.

% Um die neuen Anforderungen, die aus der CrossLab Architektur hervorgehen, zu addressieren ist eine Anpassung bzw. Neukonzeption von WIDE nötig. Dafür werden in Kapitel \ref{grundlagen} einige grundlegende Begriffe eingeführt. Danach wird in Kapitel \ref{systematische_literaturrecherche} eine Systematische Literaturrecherche durchgeführt. In Kapitel \ref{anforderungsanalyse} werden die Anforderungen an eine IDE im Kontext des GOLDi Remotelab sowie der CrossLab Architektur gesammelt. In Kapitel \ref{stand_der_technik} werden bestehende Lösungen auf ihre Anwendbarkeit überprüft und verglichen. In Kapitel \ref{konzeption} wird die Lösung konzipiert. In Kapitel \ref{implementierung} wird die Implementierung der konzipierten Lösung beschrieben. In Kapitel \ref{auswertung} wird der Erfüllungsgrad der Anforderungen ausgewertet. Schließlich wird in Kapitel \ref{zusammenfassung_und_fazit} eine Zusammenfassung der Arbeit sowie ein Fazit gegeben.
\chapter{Grundlagen}\label{section:grundlagen}

In diesem Kapitel werden einige Grundlagen für diese Arbeit beschrieben. Zunächst werden integrierte Entwicklungsumgebungen in \autoref{section:grundlagen:integrierte-entwicklungsumgebung} erläutert. Daraufhin werden das Verbundprojekt CrossLab und die dazugehörige Architektur für online Labore in \autoref{section:grundlagen:crosslab} vorgestellt.

\section{Integrierte Entwicklungsumgebung}\label{section:grundlagen:integrierte-entwicklungsumgebung}
Eine integrierte Entwicklungsumgebung bzw. \ac{IDE} ist meistens auf einen speziellen Anwendungsfall ausgelegt und besteht aus einem Code Editor sowie weiteren benötigten Softwarewerkzeugen, wie z.B. Compiler, Debugger und Language Server \cite{noauthor_language-server-protocol_nodate}. Dabei werden oftmals alle diese Komponenten direkt mit der \ac{IDE} ausgeliefert, wodurch der Nutzer direkt mit der Programmierung beginnen kann. Somit besitzt eine \ac{IDE} mehr Features als ein Code Editor, da diese nur die Bearbeitung von Code erlauben, während \acp{IDE} u.a. auch die Kompilierung und das Debuggen von Programmen ermöglichen. Language Server erweitern die Funktionen eines Code Editors, indem sie u.a. Code-Vervollständigung, Code-Navigation und Refactoring ermöglichen.

\section{CrossLab}\label{section:grundlagen:crosslab}
CrossLab \cite{aubel_adaptable_2022} ist ein Verbundprojekt der Technischen Universität Bergakademie Freiberg, der Technischen Universität Dortmund, der NORDAKADEMIE und der Technischen Universität Ilmenau. Im Rahmen des Projekts wurde eine neue Architektur für online Labore erarbeitet \cite{nau_new_2022}.

Die CrossLab-Architektur basiert auf dem Konzept von sogenannten \textit{Laborgeräten}. Diese können verschiedene \textit{Services} anbieten bzw. konsumieren. Beispiele für Services sind der \textit{Electrical Connection Service}, welcher den Austausch von Sensor- und Aktorwerten ermöglicht, und der \textit{Webcam Service}, welcher die Übertragung der Webcamaufnahmen von z.B einem elektromechanischen Hardwaremodell ermöglicht. Services besitzen immer einen \textit{Producer}, der die Funktionen des Services bereitstellt, sowie einen \textit{Consumer}, der diese nutzen kann. Dabei können auch sogenannte \textit{Prosumer} entwickelt werden, die beide Rollen erfüllen können. Durch die Verbindung der Services von verschiedenen Laborgeräten kann ein \textit{Experiment} erstellt werden. Ein Vorteil dieser Architektur ist die einfache Wiederverwendbarkeit und Austauschbarkeit von einzelnen Laborgeräten in Experimenten. So können z.B. Laborgeräte mit gleichen Services in einer entsprechenden \textit{Laborgerätegruppe} hinterlegt werden, welche dann statt eines konkreten Laborgeräts zur Erstellung eines Experiments genutzt werden kann. Beim Start des Experiments wird dann ein verfügbares Laborgerät aus der Laborgerätegruppe ausgewählt. Weiterhin gibt es noch \textit{cloud-instanziierbare} und \textit{edge-instanziierbare} Laborgeräte. Beim Start eines Experiments mit cloud- oder edge-instanziierbaren Laborgeräten wird eine entsprechende Instanz des Laborgeräts erstellt. Dabei wird für cloud-instanziierbare Laborgeräte eine Nachricht an die hinterlegte Instanziierungs-URL geschickt, wodurch die Instanz erstellt wird. Für edge-instanziierbare Laborgeräte wird mithilfe der hinterlegten Code-URL eine URL erstellt, die vom Nutzer aufgerufen werden muss um die Instanz zu erstellen.

Das Backend von CrossLab besteht aus mehreren verschiedenen Diensten, welche zusammen eine \textit{CrossLab-Instanz} bilden. Diese Dienste sind im Folgenden aufgelistet:
\begin{itemize}
    \item \textbf{Authentication Service} \\ Dieser Dienst ist für die Authentifizierung der Nutzer verantwortlich.
    \item \textbf{Authorization Service} \\ Dieser Dienst ist für die Authorisierung der Nutzer verantwortlich.
    \item \textbf{Device Service} \\ Dieser Dienst verwaltet die Laborgeräte der CrossLab-Instanz.
    \item \textbf{Experiment Service} \\ Dieser Dienst ist für die Erstellung und Verwaltung der Experimente der CrossLab-Instanz verantwortlich.
    \item \textbf{Federation Service} \\ Dieser Dienst ist für das Teilen von Laborgeräten und Experimenten mit anderen CrossLab-Instanzen verantwortlich.
\end{itemize}
Um ein Experiment in einer CrossLab-Instanz starten zu können benötigen Nutzer ein entsprechendes Nutzerkonto für diese CrossLab-Instanz.
\chapter{Systematische Literaturrecherche} \label{systematische_literaturrecherche}

\section{Vorgehensbeschreibung}

\section{Ergebnisse}

\begin{enumerate}
    \item Defining research questions
    \item Selecting databases and other research sources
    \item Defining search terms
    \item Merging hits from different databases
    \item Applying inclusion and exclusion criteria
    \item Perform the review
    \item Synthesizing results
\end{enumerate}

Um einen Überblick über den aktuellen Stand der Forschung zu bekommen wird zunächst eine systematische Literaturrecherche vorgenommen. Dabei sollen die folgenden Fragen beantwortet werden:

\begin{enumerate}
    \item Was sind Probleme von integrierten Entwicklungsumgebungen in der Lehre? /\\
          What are problems of integrated development environments in education?
    \item Was sind Vorteile von integrierten Entwicklungsumgebungen in der Lehre? /\\
          What are advantages of integrated development environments in education?
    \item Was sind Nachteile von integrierten Entwicklungsumgebungen in der Lehre? /\\
          What are disadvantages of integrated development environments in education?
    \item Was sind Anforderungen an integrierte Entwicklungsumgebungen in der Lehre? /\\
          What are requirements of integrated development environments in education?
\end{enumerate}

Die folgenden Datenbanken wurden für die Literaturrecherche ausgewählt:

\begin{itemize}
    \item ACM Digital Library
    \item DBLP Computer Science Bibliography
    \item IEEE Xplore / IEEE-IET Electronic Library Online (IEL)
    \item Web of Science
    \item zbMATH Open
\end{itemize}

Um die Datenbanken effizient durchsuchen zu können müssen nun die Suchbegriffe aus den zu beantwortenden Fragen hergeleitet werden. Dafür werden zunächst Synonyme für die einzelnen Stichwörter der Fragen gesucht. Weiterhin werden dabei auch die englischen Varianten der Fragen betrachtet, um ein möglichst breites Spektrum an möglichen Publikationen abzudecken.

\begin{table}[htbp]
    \centering
    \begin{tabularx}{\textwidth}{| >{\hsize=.6\hsize\linewidth=\hsize}X |
            >{\hsize=1.4\hsize\linewidth=\hsize}X |}
        \hline
        Stichwort                           & Synonyme                                                                                                                   \\
        \hline
        Probleme                            & Hindernisse, Hürden, Schwierigkeiten, Barrieren, Komplikationen                                                            \\ problems                            & obstacles, barriers, hurdles, impediments, obstructions, difficulties \\
        \hline
        Vorteile                            & Nutzen, positive Aspekte, Stärken                                                                                          \\ advantages                          & gains, positive aspects, benefits, strengths \\
        \hline
        Nachteile                           & negative Aspekte, Kehrseite, Schwächen                                                                                     \\ disadvantages                       & negative aspects, drawbacks, detriments, downsides, weaknesses \\
        \hline
        Anforderungen                       & Erwartungen, Ziele, Rahmenbedingungen                                                                                      \\ requirements                        & expectations, goals \\
        \hline
        integrierte Entwicklungs-umgebungen & IDEs, Code Editoren, Entwicklungsumgebungen, (Software) Entwicklungswerkzeuge, Programmierwerkzeuge, Programmierumgebungen \\ integrated development environments & IDEs, code editors, development environments, (software) development tools, programming tools, programming environments \\
        \hline
        Lehre                               & Ausbildung, Studium, Fortbildung, Weiterbildung, Universität                                                               \\ education                           & studies, training, schooling, further education, higher education, university, academic education, college, academic setting, educational setting  \\
        \hline
    \end{tabularx}
    \caption{Tabelle der Suchbegriffe}
    \label{table:search-terms}
\end{table}

Mithilfe dieser Suchbegriffe werden nun die Datenbanken durchsucht. Dafür werden zunächst für alle Fragen die entsprechenden Stichworte samt ihrer Synonyme mit OR-Operatoren verknüpft. Die auf diese Art entstandenen Gruppen für die Stichwörter werden dann für die jeweilige Frage mit AND-Operatoren verbunden. Dabei werden sowohl Singular als auch Plural der einzelnen Suchbegriffe verwendet. Die auf diese Weise gebildeten Suchanfragen werden dann in allen Datenbanken abgefragt.

\begin{table}[htbp]
    \centering
    \begin{tabular}{|l|c|c|c|c|c|c|}
        \hline
        Frage         & ACM & DBLP & IEEE & Web of Science & zbMath Open & gesamt \\
        \hline
        Probleme      &     &      &      &                &             &        \\
        \hline
        Vorteile      &     &      &      &                &             &        \\
        \hline
        Nachteile     &     &      &      &                &             &        \\
        \hline
        Anforderungen &     &      &      &                &             &        \\
        \hline
        Alle          &     &      &      &                &             &        \\
        \hline
    \end{tabular}
    \caption{Anzahl Ergebnisse der einzelnen Suchanfragen}
    \label{table:amount-search-results}
\end{table}

Die gefundenen Quellen werden nun gesammelt und analysiert. Dafür werden zunächst die obersten X Ergebnisse der Datenbanken ausgewählt und zusammengeführt. Daraufhin werden diese auf ihre Übereinstimmung mit den gestellten Forschungsfragen überprüft. Dafür werden zunächst die Überschriften, Abstracts und Keywords betrachtet. Die nach dieser Filterung verbleibenden Arbeiten werden dann zur Aufarbeitung des aktuellen Forschungsstands verwendet. Sollte die Anzahl der verbleibenden Arbeiten zu gering sein, so kann eine Vorwärts- bzw. Rückwartssuche mit Hilfe dieser erfolgen.
\chapter{Anforderungsanalyse}\label{section:anforderungsanalyse}

In diesem Kapitel werden die Anforderung an die zu entwickelnde IDE dargelegt. Diese wurden in Gesprächen mit Experten des GOLDi-Remotelab sowie Experten des CrossLab Projekts erhoben. Um einen Kontext für die Anforderungen herzustellen, wird in \autoref{section:anforderungsanalyse:beispielszenario} ein entsprechendes Beispielszenario eines Praktikumsversuchs angeführt. Daraufhin werden in \autoref{section:anforderungsanalyse:anforderungen} die erhobenen Anforderungen beschrieben.

\section{Beispielszenario: Praktikumsversuch} \label{beispielszenario}

In diesem Abschnitt soll ein Beispielszenario für den Einsatz der zu entwickelnden integrierten Entwicklungsumgebung beschrieben werden. Bei dem Beispielszenario handelt es sich um einen Praktikumsversuch. \\

Bei der Konzipierung eines Praktikumsversuchs sollten die beabsichtigten Lernergebnisse der Studierenden ermittelt werden. Anhand dieser kann dann ein entsprechender Versuchsaufbau entwickelt bzw. ausgewählt werden. Die Entwickler eines CrossLab-kompatiblen Versuchs sind dafür verantwortlich die benötigten Laborgeräte zu implementieren.

\subsection{Entwickler}

Entwickler sind für den Entwurf und die Implementierung neuer Laborgeräte verantwortlich. Gegebenenfalls muss ein neues Laborgerät mit der integrierten Entwicklungsumgebung interagieren. Um dies zu ermöglichen muss die zu entwickelnde integrierte Entwicklungsumgebung entsprechende Services anbieten. Über diese Services können dann entsprechende Laborgeräte verbunden werden. Hierbei wäre es hilfreich, wenn die integrierte Entwicklungsumgebung erweiterbar wäre, ohne das zugrundeliegende Laborgerät anpassen zu müssen.

\begin{note}
    Eine Idee wäre es in der integrierten Entwicklungsumgebung eine CrossLab-Service Erweiterungen zu erlauben. Wenn ein Experiment konfiguriert wird könnte eine entsprechende Erweiterung als Konfiguration an die integrierte Entwicklungsumgebung übergeben werden. Diese würde die entsprechende Erweiterung dann laden und könnte somit dann die darin implementierten Services anbieten. Hierbei ist allerdings noch zu überlegen, wie genau das in einem Konfigurator umgesetzt werden könnte.
\end{note}

\subsection{Betreuer}

Betreuer müssen die Lernenden während der Durchführung des Praktikumsversuchs begleiten. Hierbei könnte es dem Betreuer erlaubt werden, einem laufenden Experiment der Lernenden beizutreten. Über ein entsprechendes Rollenmanagement könnte es dem Betreuer somit ermöglicht werden den Code der Lernenden einzusehen und zu bearbeiten. Weiterhin könnte er weitere Features der integrierten Entwicklungsumgebung nutzen, wie zum Beispiel die Kompilierung und das Hochladen von Code.

\begin{note}
    Hierbei ist interessant, wie das Rollenmanagement implementiert werden muss und wie die entsprechenden Erweiterungen damit umgehen. Weiterhin ist interessant, wie das Beitreten von Nutzern allgemein implementiert werden sollte. Ideen: Betreuer kann Experiment updaten / Betreuer wird erkannt und dynamisch entsprechender Rolle zugeteilt / Optionales Laborgerät für Betreuer ist im Experiment hinterlegt. Wie sollten optionale Laborgeräte implementiert werden?
\end{note}

\subsection{Lernende}

Lernende müssen den Praktikumsversuch durchführen. Hierbei können sie alle bereitgestellten Hilfsmittel der integrierten Entwicklungsumgebung nutzen. Um die Kollaboration von Lernenden zu ermöglichen sollte die integrierte Entwicklungsumgebung über entsprechende Mechanismen verfügen. So könnten den Lernenden zum Beispiel die Kommunikation über Text und Video, sowie das gleichzeitige Editieren von Dateien ermöglicht werden.

\begin{note}
    Hierbei könnte auch über weitere Interaktionsmöglichkeiten nachgedacht werden, wie zum Beispiel eine Git-Integration oder ein Kurs-übergreifendes Wiki.
\end{note}

\subsection{Lehrende}

Lehrende sind für die didaktische Konzipierung des Praktikumsversuchs verantwortlich. Lehrende können sich darauf aufbauend ein Experiment zusammenstellen, welches ihrem didaktischen Konzept entspricht. Weiterhin könnten Lehrenden auch bereits vorkonfigurierte Experimente angezeigt werden, mit denen sie ihre angestrebten Ziele erreichen können.

\begin{note}
    Dabei ist es allerdings notwendig, dass die Geräte und Experimente mit entsprechenden Meta-Informationen ausgestattet werden, welche die Kategorisierung bzw. das Zusammenstellen von passenden Experimenten vereinfacht.
\end{note}

\subsection{Administrator}

Administratoren sind für die Verwaltung und Bereitstellung der Laborgeräte verantwortlich. Um den Administrator bei der Einrichtung der Softwarekomponenten zu Unterstützen könnte ein entsprechendes Tool konzipiert werden. Bei diesem könnte der Administrator die benötigten Softwarekomponenten auswählen und würde dann eine Konfiguration in einem passenden Format ausgegeben bekommen (z.B. Dockerfile, Docker-Compose, ...).
\section{Anforderungen}\label{section:anforderungsanalyse:anforderungen}

\newarray\Requirements
\expandarrayelementfalse
\newcounter{requirement}
\newenvironment{requirement}[1]
{
\def\reqlabel{requirement:#1}
\refstepcounter{requirement}
\label{\reqlabel}
\Requirements(\value{requirement})={\autoref{requirement:#1} & #1}
\tabularx{\textwidth}{p{4cm} X}
\toprule
\multicolumn{1}{p{3.25cm}}{\textbf{Anforderung \arabic{requirement}}} & \multicolumn{1}{c}{\textbf{#1}} \\
\midrule
}{
\bottomrule
\endtabularx
}
\newcommand{\reqdescription}{Beschreibung &}
\newcommand{\reqrationale}{Begründung &}
\def\requirementautorefname{Anforderung}
\def\replspaces#1#2{\expandafter\replspacesA\expandafter#2#1 \end}
\def\replspacesA#1#2 #3{#2\ifx\end#3\else#1\afterfi{\replspacesA#1#3}\fi}
\def\afterfi#1#2\fi{\fi#1}
\def\replace#1#2#3{%
    \def\tmp##1#2{##1#3\tmp}%
    \tmp#1\stopreplace#2\stopreplace}
\def\stopreplace#1\stopreplace{}

Im Folgenden werden die Anforderungen an die zu entwickelnde IDE beschrieben. Dabei wird für jede der Anforderungen eine entsprechende Beschreibung sowie Begründung angegeben. Eine Übersicht aller Anforderungen ist in \autoref{table:anforderungen} gegeben.

\begin{requirement}{Ausführung im Browser}
    \reqdescription Die zu entwickelnde IDE soll über den Browser des Nutzers ausgeführt werden können. \\
    \reqrationale Durch die Nutzbarkeit der IDE im Browser des Nutzers wird die Verwendung dieser vereinfacht, da Nutzer keine weitere Software installieren müssen. \\
\end{requirement}

\begin{requirement}{Erweiterbarkeit}
    \reqdescription Die zu entwickelnde IDE soll Schnittstellen zum Hinzufügen von CrossLab-Services und Benutzerinterfaces besitzen. \\
    \reqrationale Um die Weiterentwicklung der IDE zu vereinfachen sollen entsprechende Schnittstellen zur Verfügung stehen. Dabei sollte mindestens das Hinzufügen neuer CrossLab-Services und Benutzerinterfaces möglich sein. \\
\end{requirement}

\begin{requirement}{Keine Kosten für Nutzer}
    \reqdescription Die zu entwickelnde IDE soll keine Kosten für den Nutzer erzeugen. Daher soll für die Implementierung nur quelloffene und frei nutzbare Software verwendet werden.  \\
    \reqrationale Um die IDE in einer Vielzahl von verschiedenen Szenarien einsetzen zu können ist es vom Vorteil keine assoziierten Kosten für die Nutzung dieser zu haben. Somit kann sie z.B. auch im GOLDi-Remotelab und anderen CrossLab-kompatiblen online Laboren eingesetzt werden. \\
\end{requirement}

\begin{requirement}{Eigenständig nutzbar}
    \reqdescription Die zu entwickelnde IDE soll als einziges Laborgerät in einem Experiment verwendet werden können. Dabei können die angebotenen Funktionen der IDE eingeschränkt sein. Das Editieren von Quellcode soll stets ermöglicht werden. \\
    \reqrationale Da Experimente in der CrossLab-Architektur meist aus mehreren verbundenen Laborgeräten bestehen kann es dazu kommen, dass manche dieser ggf. nicht immer verfügbar sind, da sie aktuell von anderen Nutzern verwendet werden. Wenn die IDE eigenständig nutzbar ist können Nutzer dennoch an ihren Programmen weiterarbeiten. \\
\end{requirement}

\begin{requirement}{Nur CrossLab-Nutzerkonto nötig}
    \reqdescription Die zu entwickelnde IDE soll nur das bereits zur Ausführung von CrossLab Experimenten benötigte Nutzerkonto voraussetzen. \\
    \reqrationale Durch die Notwendigkeit eines zweiten Nutzerkontos könnte die IDE für gewisse Nutzergruppen uninteressant werden. Somit soll nur ein CrossLab-Nutzerkonto für die Nutzung der IDE vorausgesetzt werden. \\
\end{requirement}

\begin{requirement}{Integriertes Dateisystem}
    \reqdescription Die zu entwickelnde IDE soll ein integriertes Dateisystem besitzen. \\
    Unteranforderung (a) & Das integrierte Dateisystem soll die Erstellung, Bearbeitung, Verschiebung, Löschung und persistente Speicherung von Dateien und Ordnern unterstützen. \\
    Unteranforderung (b) & Das integrierte Dateisystem der IDE soll bei der eigenständigen Ausführung unterstützt werden. (sh. \autoref{requirement:Eigenständig nutzbar}) \\
    \reqrationale Ein in der IDE integriertes Dateisystem vereinfacht die Nutzung der IDE, da kein externes Dateisystem benötigt wird. Durch die Unterstützung des integrierten Dateisystems bei der eigenständigen Ausführung der IDE haben Nutzer zudem weiterhin Zugriff auf ihre Dateien. \\
\end{requirement}

\begin{requirement}{Weitere Dateisysteme}
    \reqdescription Die zu entwickelnde IDE soll die Anbindung weiterer Dateisysteme unterstützen. \\
    Unteranforderung (a) & Es sollen CrossLab-Services für die Bereitstellung und Nutzung von Dateisystemen entwickelt werden. \\
    Unteranforderung (b) & Die zu entwickelnde IDE soll die aus Unteranforderung (a) resultierenden CrossLab-Services zur Anbindung weiterer Dateisysteme nutzen. \\
\end{requirement}

\begin{requirement}{Kollaboration}
    \reqdescription Die zu entwickelnde IDE soll Echtzeit-Kollaboration durch die Synchronisation von geteilten Daten und den Austausch von Zustandsinformationen ermöglichen. \\
    Unteranforderung (a) & Es sollen Crosslab-Services für die Synchronisation von geteilten Daten und den Austausch von Zustandsinformationen entwickelt werden. \\
    Unteranforderung (b) & Die zu entwickelnde IDE soll die aus Unteranforderung (a) resultierenden CrossLab-Services zur Bereitstellung der Echtzeit-Kollaboration verwenden. \\
    \reqrationale Durch die Ermöglichung der Synchronisation von Daten und dem Austausch von Zustandsinformationen zwischen Nutzern innerhalb eines Experiments kann die Zusammenarbeit dieser gefördert werden. \\
\end{requirement}

\vfill

\begin{requirement}{Teilen von Ordnern}
    \reqdescription Die zu entwickelnde IDE soll das Teilen von Ordnern und den darin enthaltenen Dateien zwischen Nutzern innerhalb eines Experiments ermöglichen. \\
    Unteranforderung (a) & Änderungen innerhalb geteilter Ordner sollen zwischen allen teilnehmenden Nutzern synchronisiert werden. \\
    Unteranforderung (b) & Geteilte Ordner können nur von ihrem Besitzer gelöscht, verschoben oder umbenannt werden. \\
    Unteranforderung (c) & Das Teilen von Ordnern soll von dem Besitzer beendet werden können. \\
    Unteranforderung (d) & Das Teilen von Ordnern soll über die nach \autoref{requirement:Kollaboration}(b) vorhandenen CrossLab-Services der IDE erfolgen. \\
    \reqrationale Durch das Teilen von Ordnern und den enthaltenen Dateien können Nutzer gemeinsam an diesen arbeiten. Dadurch können z.B. Gruppenarbeiten im Rahmen eines Praktikumsversuch effizienter durchgeführt werden, während die Lernenden gleichzeitig ihre Teamfähigkeit verbessern können. \\
\end{requirement}

\newpage

\mbox{}\vfill

\begin{requirement}{Kompilierung}
    \reqdescription Die zu entwickelnde IDE soll die Kompilierung von Quellcode unterstützen. \\
    Unteranforderung (a) & Es sollen CrossLab-Services für die Bereitstellung und Nutzung von Compilern entwickelt werden. \\
    Unteranforderung (b) & Die zu entwickelnde IDE soll die aus Unteranforderung (a) resultierenden CrossLab-Services für die Nutzung von Compilern verwenden. \\
    \reqrationale Die Kompilierung des Quellcodes ist in vielen Programmiersprachen ein wichtiger Schritt um das Programm auf dem Zielsystem ausführen zu können. \\
\end{requirement}

\vfill

\begin{requirement}{Programmierung von Steuereinheiten}
    \reqdescription Die zu entwickelnde IDE soll die Programmierung Steuereinheiten unterstützen. \\
    Unteranforderung (a) & Es soll ein CrossLab-Service für Programmierung von Steuereinheiten entwickelt werden. \\
    Unteranforderung (b) & Die zu entwickelnde IDE soll den aus Unteranforderung (a) resultierenden CrossLab-Service für die Programmierung von Steuereinheiten verwenden. \\
    \reqrationale Nutzer können innerhalb eines Experiments Programme für Steuereinheiten schreiben. Zur Ausführung müssen die Programme auf die entsprechende Steuereinheit hochgeladen und für deren Programmierung verwendet werden. \\
\end{requirement}

\vfill\mbox{}

\newpage

\begin{requirement}{Debuggen}
    \reqdescription Die zu entwickelnde IDE soll das Debuggen von Programmen der Nutzer ermöglichen. \\
    Unteranforderung (a) & Es soll ein CrossLab-Service für die Bereitstellung und Nutzung von Debuggern entwickelt werden. \\
    Unteranforderung (b) & Es soll ein CrossLab-Service für die Kommunikation zwischen Debuggern und Steuereinheit entwickelt werden. \\
    Unteranforderung (c) & Die zu entwickelnde IDE soll den aus Unteranforderung (a) resultierenden CrossLab-Service für die Nutzung von Debuggern verwenden. \\
    \reqrationale Nutzer können durch das Debuggen ihrer Programme schneller Fehler in diesen finden und beheben. Weiterhin erlaubt das Debuggen eines Programms einen besseren Einblick in dessen Laufzeitverhalten. \\
\end{requirement}

\vfill

\begin{requirement}{Teilen von Debug-Sitzungen}
    \reqdescription Die zu entwickelnde IDE soll es Nutzern ermöglichen gleichzeitig an einer Debug-Sitzung teilzunehmen, falls dies mit dem verwendeten Debugger möglich ist. \\
    Unteranforderung (a) & Nutzer können einer Debug-Sitzung nur beitreten wenn sie Zugriff auf die Dateien des entsprechenden Programms besitzen. (sh. \autoref{requirement:Teilen von Ordnern}) \\
    Unteranforderung (b) & Die Breakpoints aller an einer Debug-Sitzung teilnehmenden Nutzer sollen synchronisiert werden. \\
    Unteranforderung (c) & Pro Steuereinheit soll nur eine Debug-Sitzung gestartet werden können. \\
    Unteranforderung (d) & Nur der Ersteller einer Debug-Sitzung kann diese beenden. \\
    Unteranforderung (e) & Das Teilen von Debug-Sitzungen soll über den nach \autoref{requirement:Kollaboration}(b) vorhandenen CrossLab-Service der IDE erfolgen. \\
    \reqrationale Durch das kollaborative Debuggen können Nutzer gemeinsam einen Einblick in das Laufzeitverhalten des Programs erlangen. Dadurch kann auch die Suche nach Fehlern sowie deren Behebung effizienter erfolgen. \\
\end{requirement}

\newpage

\begin{requirement}{Testen}
    \reqdescription Es soll ein CrossLab-Service für die Erstellung und Ausführung von Testfällen innerhalb eines Experiments entwickelt werden. \\
    Unteranforderung (a) & Der Producer des CrossLab-Service soll die Bereitstellung von Funktionen zur Verwendung in Testfällen ermöglichen. \\
    Unteranforderung (b) & Der Consumer des CrossLab-Service soll die Ausführung der bereitgestellten Funktionen ermöglichen. \\
    Unteranforderung (c) & Die Erstellung von Testfällen soll während der Erstellung eines Experiments erfolgen können. \\
    Unteranforderung (d) & Die zu entwickelnde IDE soll den CrossLab-Service zur Ausführung von Testfällen innerhalb eines Experiments verwenden. \\
    \reqrationale Die Möglichkeit Testfälle für ein Experiment zu konfigurieren erlaubt es z.B. Lehrenden die Ziele ihrer Lehrveranstaltung im Vorhinein festzulegen. Während des Experiments können die Lernenden dann ihre erstellte Lösung überprüfen. \\
\end{requirement}

\vfill

\begin{requirement}{Language Server}
    \reqdescription Die zu entwickelnde IDE soll die Anbindung von Language Servern unterstützen. \\
    Unteranforderung (a) & Es soll ein CrossLab-Service für die Bereitstellung und Nutzung von Language Servern entwickelt werden. \\
    Unteranforderung (b) & Die zu entwickelnde IDE soll den aus Unteranforderung (a) resultierenden CrossLab-Service für die Nutzung von Language Servern verwenden. \\
    \reqrationale Durch die Anbindung von Language Servern können Editorfunktionen wie Code-Vervollständigung, Code-Navigation und Refactoring ermöglicht werden. Diese können die Benutzererfahrung verbessern. \\
\end{requirement}

\newpage

\begin{table}[t]
    \centering
    \begin{tabular}{l l}
        \toprule
        \Requirements(1)  \\
        \Requirements(2)  \\
        \Requirements(3)  \\
        \Requirements(4)  \\
        \Requirements(5)  \\
        \Requirements(6)  \\
        \Requirements(7)  \\
        \Requirements(8)  \\
        \Requirements(9)  \\
        \Requirements(10) \\
        \Requirements(11) \\
        \Requirements(12) \\
        \Requirements(13) \\
        \Requirements(14) \\
        \Requirements(15) \\
        \bottomrule
    \end{tabular}
    \caption{Übersicht der Anforderungen}
    \label{table:anforderungen}
\end{table}

\chapter{Stand der Technik} \label{stand_der_technik}

\begin{itemize}
    \item Standalone IDEs / Code Editoren
          \begin{itemize}
              \item VSCode
              \item Brackets
              \item Phoenix
              \item IntelliJ IDEs
          \end{itemize}
    \item IDE Frameworks
          \begin{itemize}
              \item Theia
              \item OpenSumi
          \end{itemize}
    \item Remote Development Platforms
          \begin{itemize}
              \item Eclipse Che
              \item Coder
              \item Gitpod
              \item Replit
              \item Codeanywhere
              \item Stackblitz
              \item Github Codespaces
              \item Amazon Cloud9
          \end{itemize}
    \item Educational
          \begin{itemize}
              \item Codeboard
              \item Coding Rooms
              \item CodeHS
              \item JDoodle
              \item KIRA
              \item Codio
              \item Online IDE
              \item Codegrade
          \end{itemize}
\end{itemize}

Im Bereich der IDEs gibt es bereits viele bestehende Softwarelösungen. Für die Zwecke dieser Arbeit ergeben sich drei Kategorien:

\begin{itemize}
    \item \textbf{Standalone IDEs / IDE-Frameworks:} \\ Diese Kategorie beinhaltet Softwarelösungen wie VSCode, Theia und OpenSumi.
    \item \textbf{Educational IDE Plattformen:} \\ Diese Kategorie beinhaltet Softwarelösungen wie CodeHS, Codegrade und KIRA.
    \item \textbf{Workspace Management:} \\ Diese Kategorie beinhaltet Softwarelösungen wie Eclipse Che, Coder und Gitpod.
\end{itemize}

Im folgenden werden die einzelnen Kategorien genauer betrachtet.

\section{Standalone IDEs / IDE-Frameworks}
In der Kategorie Standalone IDEs / IDE-Frameworks befinden sich unter anderem Visual Studio Code, Theia und OpenSumi.

\section{Educational IDE Plattormen}


\section{Workspace Management}
\chapter{Konzeption} \label{konzeption}

Im Fokus dieser Arbeit liegt die Programmierung von Mikrocontrollern im Rahmen des GOLDi Remotelab. Bei den verwendeten Microcontrollern handelt es sich um ATmega2560. Damit diese mit der CrossLab Infrastruktur kommunizieren können sind sie über einen FPGA mit einem Raspberry Pi Compute Module 4 verbunden. Dabei übernimmt der FPGA die Kommunikation zwischen dem CM und dem Microcontroller, während das CM die Kommunikationsschnittstelle zur CrossLab Infrastruktur übernimmt. Als Beispiel für ein steuerbares elektromechanisches Modell wird das 3-Achs-Portal verwendet. Neben den realen Systemen sollen auch die virtuellen Versionen in Betracht gezogen werden. Daraus folgen vier verschiedene Experiment-Konfigurationen, die für die Konzeption betrachtet werden:

\begin{enumerate}
    \item Realer Microcontroller und reales 3-Achs-Portal
    \item Realer Microcontroller und virtuelles 3-Achs-Portal
    \item Virtueller Microcontroller und reales 3-Achs-Portal
    \item Virtueller Microcontroller und virtuelles 3-Achs-Portal
\end{enumerate}

Weiterhin ist zu beachten, dass auch mehrere Steuereinheiten und Modelle in einem Experiment enthalten sein können.

\section{Datenspeicherung} \label{datenspeicherung}

Um sicherzustelllen, dass die IDE komplett im Browser verwendet werden kann, müssen die Nutzer in der Lage sein ihre Projekte in dieser verwalten zu können. Da noch nicht alle Browser die File System API komplett implementieren muss die Projektverwaltung auf einem anderen Weg implementiert werden. Dazu kann die IndexedDB API genutzt werden. Diese ermöglicht die persistente Speicherung größerer Datenmengen direkt im Browser. Um die IndexedDB zur Projektverwaltung nutzen zu können muss eine entsprechende Extension für Visual Studio Code geschrieben werden. Dafür kann die FileSystemProvider API verwendet werden. Diese ermöglicht es eigene Dateisysteme in Visual Studio Code einzubinden. Weiterhin gibt es noch die FileSearchProvider API und die TextSearchProvider API. Die FileSearchProvider API ermöglicht es die Logik für die Suche nach Dateien in einem bestimmten Dateisystem zu implementieren. Die TextSearchProvider API ermöglicht es die Logik für die Durchsuchung des Textinhalts von Dateien zu implementieren. Sowohl die FileSearchProvider API als auch die TextSearchProvider API befinden sich aktuell noch in einem experimentellen Zustand. Das bedeutet, dass sie verwendet werden können, allerdings dürfen Extensions, die diese APIs nutzen nicht auf dem Marketplace veröffentlicht werden. Dies sollte allerdings für die Verwendung innerhalb des GOLDi Remotelab kein Problem darstellen.

Weiterhin besteht jedoch auch die Möglichkeit die Projekte der Nutzer auf einem Server zu speichern. Dies würde es den Nutzern ermöglichen zwischen ihren Endgeräten wechseln zu können, ohne ihren aktuellen Stand zu verlieren. Allerdings müssen dann entsprechende Server bereitgestellt werden. Ein anderer Ansatz könnte es den Nutzern erlauben ihre Projekt aus GitHub oder GitLab zu laden.
\section{Kompilierung} \label{kompilierung}

Die Kompilierung des Quellcodes kann vom CM übernommen werden. Alternativ steht auch hier die Möglichkeit eines externen Servers zur Verfügung. Weiterhin kann auch die Möglichkeit der Kompilierung innerhalb des Browsers des Nutzers erforscht werden. Hierbei gibt es bereits eine Version des C-Compilers clang in WebAssembly. Diese ermöglicht jedoch in der aktuellen Form nur die Kompilierung von C/C++ Code zu WebAssembly. Im Falle des GOLDi Remotelab müsste allerdings eine Kompilierung des Quellcodes für die verwendeten Microcontroller erfolgen. Diese wird normalerweise mit Hilfe des Compilers avr-gcc durchgeführt. Allerdings bietet auch der clang Compiler ein avr Ziel an, welches jedoch einen experimentellen Status hat.
\section{Language Server} \label{language_server}

\begin{itemize}
    \item Language Server über Cloud-instanziierbares Gerät oder über clangd im Browser
\end{itemize}

Eine Möglichkeit zur Einbindung von Language Servern ist es, diese auf dem CM zu installieren und sich dann zu diesen zu verbinden. Allerdings ist das Language Server Protokoll nur auf einen einzigen Client pro Server ausgelegt. Daher müsste in einem Experiment mit mehreren Nutzern für jeden dieser Nutzer eine eigene Instanz des Language Server erstellt werden. Eine spezielle Alternative könnte die Verwendung des zu WebAssembly kompilerten Language Servers clangd darstellen. Dieser könnte im Browser des jeweiligen Nutzers instanziiert werden, wodurch keine zusätzliche Last auf dem CM entsteht. Allerdings muss hierbei die neu entstandene Last auf dem Rechner des Nutzers in Betracht gezogen werden. Eine weitere Möglichkeit ist die Bereitstellung weiterer Server zur Bereitstellung der Language Server.
\section{Debugging} \label{debugging}

\begin{itemize}
    \item Debugging von realem Microcontroller über RPi (avr-gcc)
    \item Debugging von virtuellem Microcontroller über Cloud-instanziierbares Gerät oder über avr-gcc im Browser
\end{itemize}

Das Debuggen von Programmen auf den vorhandenen Microcontrollern gestaltet sich schwierig. Eine Möglichkeit ist die Nutzung der Bibliothek avr\_debug. Diese wird zusammen mit dem Programm kompiliert und auf den Microcontroller hochgeladen. Dort erstellt sie ein Interface für den Debugger gdb. Dieses Interface nutzt die Serielle Schnittstelle des Microcontrollers zur Kommunikation mit gdb. Das CM agiert in diesem Szenario als Schnittstelle zwischen gdb und unserer IDE. Ein Nachteil dieses Vorgehens ist der hohe Speicherverbrauch der Bibliothek, welcher die Anzahl möglicher Programme einschränkt. Allerdings ist ein Vorteil dieses Ansatzes, dass keine zusätzlichen Kosten durch die Anschaffung externe Debugger entstehen.

Ein weiteres Problem, was beim Debuggen eines laufenden Experimentes beachtet werden muss, ist die fortlaufende Ansteuerung von weiteren Geräten. Nehmen wir als Beispiel ein einfaches Experiment bestehend aus einem Microcontroller und einem 3-Achs-Portal. Wenn wir das Program des Microcontrollers unterbrechen, während dieser den Portalkran aktiv nach rechts bewegt, so wird diese Bewegung nicht unterbrochen. Um sicherzustellen, dass die Signale von Aktoren während eines Breakpoints nicht an andere Geräte weitergeleitet werden müssen die anderen Geräte entsprechend benachrichtigt werden.
\section{Kollaboratives Zusammenarbeiten} \label{kollaboration}

\begin{itemize}
    \item Ein Nutzer teilt sein Projekt mit den anderen Nutzern
    \item Alle Nutzer greifen auf ein geteiltes Projekt zu
    \item Kompilierung, Upload und Debugging sollen auch geteilt werden
\end{itemize}

Um den Nutzern des GOLDi Remotelab ein kollaboratives Zusammenarbeiten zu ermöglichen muss eine entsprechende Extension konzipiert werden. Dazu ist es zunächst wichtig, unterschiedliche Ansätze von Kollaboration zu vergleichen. Durch den Vergleich soll ermittelt werden, welche Ansätze sich für den Einsatz im GOLDi Remotelab eignen. Dabei werden die folgenden Ansätze betrachtet:

\begin{enumerate}
    \item Ein Nutzer teilt sein Projekt mit anderen Nutzern innerhalb eines Experiments
    \item Nutzer arbeiten zusammen an einem geteilten Projekt
\end{enumerate}

\subsection{Ansatz 1}

In Ansatz 1 besitzt nur ein Nutzer das Projekt. Das bedeutet, dass nur bei diesem Nutzer der aktuelle Stand des Projekts gespeichert wird. Bei allen anderen Nutzern werden lediglich die Änderungen während eines Experiments synchronisiert. Sobald das Experiment beendet wurde werden keine weiteren Änderungen am Projekt mehr synchronisiert.

\subsection{Ansatz 2}

In Ansatz 2 wird der aktuelle Zustand des Projekts auf einem Server gespeichert. Alle Nutzer mit den entsprechenden Berechtigungen können dementsprechend auf den aktuellen Stand des Projekts zugreifen und Änderungen vornehmen, solange sie eine aktive Verbindung zum Server haben. Sollten sie keine aktive Verbindung zum Server haben, sollten Ihnen keine Änderungen am Projekt erlaubt sein.
\chapter{Implementierung} \label{implementierung}

\section{CrossLab Kompatibilität}

\begin{itemize}
    \item Definition und Implementierung der Services
    \item Anbindung an IDE
\end{itemize}

\section{Datenspeicherung}

\begin{itemize}
    \item FileSystemProvider
    \item FileSearchProvider (Proposed API)
    \item TextSearchProvider (Proposed API)
    \item Speicherung in indexeddb an Domain gebunden $\to$
\end{itemize}

\section{Code Kompilierung}

\begin{itemize}
    \item Kompilierungsserver
    \item Kompilierung auf Raspberry Pi
\end{itemize}

\section{Language Server}

\begin{itemize}
    \item Server wahrscheinlich benötigt
    \item clangd WebAssembly? (\href{https://github.com/Guyutongxue/clangd-in-browser}{https://github.com/Guyutongxue/clangd-in-browser})
    \item Handhabung kollaboratives Arbeiten an Code
\end{itemize}

\section{Debugging}

\begin{itemize}
    \item avr\_debug (\href{https://github.com/jdolinay/avr_debug}{https://github.com/jdolinay/avr\_debug})
    \item Anbindung des Debuggers an das Frontend
\end{itemize}

\section{Kollaboratives Bearbeiten}

\begin{itemize}
    \item Yjs kann genutzt werden (\href{https://github.com/yjs/yjs}{https://github.com/yjs/yjs})
    \item wahrscheinlich Anpassung des webrtc-connectors erforderlich \\ (\href{https://github.com/yjs/y-webrtc}{https://github.com/yjs/y-webrtc})
\end{itemize}
\chapter{Auswertung} \label{auswertung}
\chapter{Zusammenfassung und Fazit} \label{zusammenfassung_und_fazit}

% opt: zusätzliche Anhänge
\begin{appendix}
	% \input{content/anhang1}
\end{appendix}

% --- Abbildungs-, Tabellen-, Algorithmen- und Literaturverzeichnis ------------
\backmatter
% Abbildungsverzeichnis
\listoffigures
% Tabellenverzeichnis
\listoftables
% Algorithmenverzeichnis
\listofalgorithms
% Literaturverzeichnis
\cleardoublepage	% Literaturverzeichnis immer auf ungerader Seite
\phantomsection		% Anker für Sprungmarke im Inhaltsverzeichnis korrigieren
\cite{badrinarayanan2015segnet}
\addcontentsline{toc}{chapter}{\bibname}
\bibliography{NIKR_bibliography}
\end{document}
